
\subsection{Tutorial}%
\label{sub:Tutorial}
The tutorial serves to orient the user, introduce the voting rules,
and give the user a sense of how to interact with the application.

It begins by defining the spatial model, contextualizing the axes
as representing a policy or ideology space, and emphasizing the
key assumption: only the distance between voters and candidates matters
in determining how voters will fill out their ballots.

Next, it introduces the tool and the primary ways of using it
via a series of examples that introduce Plurality, Borda, IRV, and STV.
At each step, the tutorial shows an example plot
and asks the user to make a prediction of which candidate will win.

The plurality example demonstrates the well-known "spoiler effect",
in which a majority bloc of voters is shut out of representation
because they split their votes across multiple candidates,
while the minority bloc is unified in their support of a single candidate.

The Borda example explains how candidates are scored,
then demonstrates its tendency to elect centrist "consensus" candidates.

The IRV example shows how runoff voting can avoid the spoiler effect:
the plot is set up the same as in the plurality example,
but, in this case, one of the candidates supported by the majority bloc
wins the election.

Finally, the multi-winner example contrasts bloc plurality with STV.
It shows two clusters of voters, one slightly larger than the other,
with two candidates per voter cluster.
In bloc plurality, the larger cluster can sweep the election,
securing both seats on the committee, shutting out the other cluster entirely.
However, in STV, each cluster gets one seat on the committee---a more proportional outcome.

\subsection{Input}%
\label{sub:Input}
On the plot, voters are blue dots, and candidates are yellow stars.
There are four primary methods for placing voters and candidates on the plot.
\begin{itemize}
  \item \textbf{Manual Input.}
    The user can select to add a candidate or a voter,
    then click anywhere in the plot to add a point at that location.

  \item \textbf{Synthetic Sampling.}
    In the sidebar, the user can choose from any of the available distributions.
    They can set the distribution center by dragging it on the plot,
    or specifying the desired coordinate position in an input box.
    They can set distribution parameters with sliders or by entering
    the desired value in an input box.
    Then, they can set whether to sample voters or candidates,
    and set the number of points to be sampled.
    Finally, clicking "add data" adds these points to the plot.

    By stacking sets of samples drawn from different distributions,
    the user can easily create multiple clusters of voters or candidates.

    As in \cite{elkind}, the supported distributions are
    the Gaussian, the Uniform Rectangle, and the Uniform Disc.

  \item \textbf{Survey Data.}
    Upon selecting the "Survey" tab in the sidebar,
    the user can select two "Ideology" axes or political issues
    to serve as the axes of the plot.
    Then, they can sample voters from polling data.

    \item \textbf{Data Upload.}
    If a user has exported a JSON file of the points from a previously-generated plot,
    they may upload that data.
\end{itemize}

\subsection{Plot Management}%
\label{sub:Plot Management}
Every point manually added and every distribution sampled from
appears as a separate line item in the "Data Points" section of the sidebar.
The line describes the parameters that created it.
For every distribution line, the following interactions are available:
\begin{itemize}
  \item \textbf{Edit.}
    Highlights the distribution and reveals its settings;
    any of the settings (include type, count, position, and parameters)
    can be changed, and when the user clicks "update",
    the points are resampled.
    The user can click "cancel" at any point in this interaction
    if they want to keep the old settings.
  \item \textbf{Delete.}
    Removes this distribution.
  \item \textbf{Hide.}
    Hides this distribution; it is temporarily disabled
    for any election runs, but stays listed in the sidebar
    to be easily re-enabled later.
\end{itemize}
For any line item representing a manually-placed candidate or voter,
only the Delete and Hide options are available.

\subsection{Running Elections}%
\label{sub:Running Elections}
At the bottom of the sidebar, the user can choose their desired voting method,
select the number of winners, and run the election.
Running the election updates the plot,
making the winning candidates visually distinct from the others.


\subsection{Survey Data}%
\label{sub:Survey Data}

\subsection{Implementation}%
\label{sub:Implementation}
The application front end is built with Svelte and D3 in JS,
with Vite as the server.
The election algorithms are implemented in Python and
compiled to Web Assembly via Pyodide so that they can be run in browser.

